\documentclass[10pt,a4paper]{amsart}
\usepackage[margin=19mm]{geometry}
\usepackage{amsmath,amssymb,mathtools}
\usepackage[scr=rsfs,cal=euler]{mathalfa}
\usepackage{xfrac}
\usepackage[unicode,hidelinks,bookmarks=false]{hyperref}
\newcommand{\ie}{i.e.\ }
\newcommand{\cf}{cf.\ }
\newcommand{\resp}{respectively\ }
\newcommand{\Subsection}{\subsection}
\providecommand{\nobreakdash}{\nobreak\hbox{-}}
\setlength{\emergencystretch}{2em}
\multlinegap0pt
\usepackage{bm}
\usepackage{bbm}
\usepackage{dsfont}
\usepackage{xspace}
\usepackage{cancel}
\usepackage{mathtools}
\usepackage{amsthm}
\makeatletter
\@ifundefined{thm}{\newtheorem{thm}{Theorem}}{}
\@ifundefined{cor}{\newtheorem{cor}[thm]{Corollary}}{}
\@ifundefined{lemma}{\newtheorem{lemma}[thm]{Lemma}}{}
\@ifundefined{prop}{\newtheorem{prop}[thm]{Proposition}}{}
\makeatother
\title[Disintegrating the curve complex]{Disintegrating the curve complex}
\author{Mladen Bestvina, Kenneth Bromberg, Alexander J. Rasmussen}
\date{Source: arXiv:2510.03980v1 (CC BY 4.0)}
\begin{document}
\maketitle

\noindent\emph{Source and licence.} Disintegrating the curve complex. Version arXiv:2510.03980v1 (CC BY 4.0), distributed under the Creative Commons Attribution 4.0 International licence, as recorded in the arXiv licence metadata for this exact version and shown on the abstract page https://arxiv.org/abs/2510.03980v1. Full rights notice: licenses/arxiv-2510.03980-CC-BY-4.0.txt.\par\smallskip

\noindent\textit{Editorial scope.} Counted unit: the Thm whose source label is \texttt{quasitreelem} in the article (source0.tex lines 700--709, source label \texttt{quasitreelem}), delivered together with the article's own complete original proof of it (source0.tex lines 710--765, one \texttt{proof} environment of 56 lines). No mathematics is written, repaired or re-derived: statement and proof are the article's own text line for line. Required context retained uncounted: fellow travel (lines 579--596); intersection quasiconvex (lines 583--585); larger (lines 583--585); triple (lines 583--585); union quasiconvex (lines 583--585); geodesic\_bottleneck (lines 637--640); bottleneck\_converse (lines 643--649); subtrees (lines 678--691). Every definition, display and label of those lines is retained unchanged. Omitted from this package and not redistributed: 1--578, 586--636, 641--642, 650--677, 692--699, 766--2757. Nothing inside a retained excerpt is omitted or altered: every retained body is delivered exactly as the article's lines are written, so no omission receipt of any kind accompanies this package. Citations: the retained excerpts carry no citation command at all, so no bibliography entry of the article is transcribed and no reference is redistributed. Reference adaptation: none was needed and none was made; every reference inside the delivered unit resolves inside it, so no unresolved reference or citation is produced. Printed slips kept literal: no printed word, symbol, hypothesis or number of the counted statement or of its proof was altered. Layout and class: the article's own class is \texttt{article} with packages including latexsym, comment, amsmath, amsthm, amsfonts, amssymb, graphicx, epsfig, color, epsf, enumerate, tikz-cd, mathtools, float; the wrapper is the lane's pinned \texttt{amsart} editorial wrapper at 10pt on A4, which restates the theorem-like environments the retained text needs and adds the article's own macro definitions that the retained text needs (none), with extra wrapper packages (bm, bbm, dsfont, xspace, cancel, mathtools). The delivered numbering is the unit's own. Assets: no retained span contains a figure environment, a graphics-inclusion command, a PDF-page inclusion, a table or a TikZ picture; no image file is copied into this package, the package declares no asset dependency and no image file is redistributed with it. Licence: version arXiv:2510.03980v1 of this article is distributed under the Creative Commons Attribution 4.0 International licence, as recorded in the arXiv licence metadata for this exact version and shown on the abstract page https://arxiv.org/abs/2510.03980v1. A full rights notice is delivered at licenses/arxiv-2510.03980-CC-BY-4.0.txt. Characters: every retained excerpt is pure ASCII, so the delivered engine, XeLaTeX with its default Latin Modern text font, reports no missing character.
\begin{lemma}\label{fellow travel}
Fix $0<\delta\le Q \le C$ and let $X$ be a $\delta$-hyperbolic geodesic metric space.
\begin{enumerate}
\item Let $A$ be a $Q$-quasi-convex subset of $X$. If $x,y \in B_C(A)$ then 
\begin{eqnarray*}
[x,y] &\subset & B_{Q+2\delta}(A) \cup B_{C+2\delta}(\{x,y\})\\  & \subset & B_{C+2\delta}(B_C(A)).
\end{eqnarray*}

\item\label{intersection quasiconvex} If $A_1, \dots, A_n$ are $Q$-quasi-convex subsets of $X$  then $I_C(A_1, \dots, A_n)$ is $(C+2\delta)$-quasi-convex.

\item\label{union quasiconvex} If $A_0, A_1, \dots$ are $Q$-quasi-convex subsets of $X$ with $I_C(A_0, A_i) \neq \emptyset$ for all $i\ge 1$ and $C\ge Q$, then $\bigcup_{i\ge 0} A_i$ is $(2C+4\delta)$-quasi-convex.

\item\label{triple} If $A_1, A_2$, and $A_3$ are $Q$-quasi-convex subsets and $I_C(A_1, A_2)$, $I_C(A_1, A_3)$, and $I_C(A_2, A_3)$ are all non-empty, then $I_{2C+3\delta}(A_1, A_2, A_3) \neq \emptyset$.

\item\label{larger}  If $A_1, \dots, A_n$ are $Q$-quasi-convex subsets of $X$ and the diameter of $I_C(A_1, \dots, A_n)$ is $> 4C + 8\delta$, then
$$I_C(A_1, \dots, A_n) \subset B_{C+2\delta}(I_{Q+2\delta}(A_1, \dots, A_n)).$$
\end{enumerate}
\end{lemma}

\begin{cor}\label{geodesic_bottleneck}
Let $A$ be a $Q$-quasi-convex subgraph of a graph $X$ and let $K \ge 2Q+1$ and $\Delta>0$. If $A$ satisfies a
$(\Delta, K)$-bottleneck condition then $A$ is quasi-isometric to a simplicial tree with quasi-isometry constants depending only on $K$ and $\Delta$.
\end{cor}

\begin{prop}\label{bottleneck_converse}
Let $A$ be a $Q$-quasi-convex subgraph of a $\delta$-hyperbolic graph $X$. If $A$
(with the subspace metric) is quasi-isometric to a simplicial tree then,
for any $K\ge 2Q+1$, there is $\Delta>0$, depending on $K$ and the quasi-isometry constants, such that $A$ satisfies a $(\Delta, K)$-bottleneck condition.

Furthermore, for $D>0$ any $D$-connected subset of $A$ is a quasi-tree.
\end{prop}

\begin{lemma}\label{subtrees}
Let $T_0$ and  $T_1$ be $Q$-quasi-convex subsets of a $\delta$-hyperbolic graph that are quasi-trees with $I_C(T_0, T_1) \neq \emptyset$. Then there are quasi-trees $T'_0$ and $T^i_1$ with 
 $T_0 \cup T_1 = T'_0 \cup \left(\bigcup_j T^i_1\right)$
where
\begin{itemize}
\item $T_0 \subset T'_0\subset B_{C'}(T_0)$;

\item for each $i$, $T_1^i$ is $Q^*$-quasi-convex and $T_1^i \subset T_1$;

\item for each $i$, $I_{C'}(T_0, T^i_1) \neq \emptyset$;

\item for $0\le i<j$, the diameter of $I_{Q^*+2\delta}(T^i_1, T^j_1)$ is $< D$.
\end{itemize}
\end{lemma}

\begin{thm}\label{quasitreelem}
Let $X$ be a $\delta$-hyperbolic graph and let $T_0, T_1,\dots$ be subgraphs that are $Q$-quasi-convex, uniform quality quasi-trees.
Assume that
\begin{enumerate}
\item there exists $C>0$ such that $I_{C}(T_0, T_i) \neq \emptyset$ for all
$i\ge 1$;
\item there exists $N>0$ and $D>0$ such that the diameter of $I_{Q+2\delta}(T_{i_1}, \dots, T_{i_N})$ is bounded by $D$ for all $1\le i_1 < i_2 < \cdots < i_N$.
\end{enumerate}
Then $T = \displaystyle{\bigcup_{i\ge 0}} T_i$ is a quasi-tree with constants only depending on $\delta$, $Q$, and the quasi-tree constants for the sets $T_i$.
\end{thm}

\begin{proof}
By Lemma \ref{fellow travel} part \eqref{union quasiconvex}, $T$ is quasi-convex. We induct on $N$ to prove that $T$ is a quasi-tree.

The base case is $N=1$. We first prove this with the extra assumption that the diameter of $I_{Q+2\delta}(T_i, T_j)$ is also bounded above by $D$ for all
$i,j$.


By Lemma \ref{fellow travel} part \eqref{union quasiconvex}, $T$ is $(2C+4\delta)$-quasi-convex and is $K$-connected with $K = 4C + 8\delta +1$.
Furthermore, if $I_K(T_i, T_j) \neq \emptyset$ with $0<i<j$ then, by Lemma \ref{fellow travel} part \eqref{triple}, we have that $I_{K'}(T_0, T_i, T_j) \neq \emptyset$ with $K' = 2K+3\delta$ and, using our extra assumption, by \eqref{larger} the diameter of $I_{K'}(T_0, T_i, T_j)$ is $< D'$ where $D' = 2K'+4\delta + D$.


For each $T_i$ with $i\ge1$, choose $t_i \in  I_{C}(T_0, T_i)$. If $x \in T_i$ and $y \in T_j$ with $0\le i<j$ and $d(x,y) \le K$ then $d(x, t_i) \le 2D'$. Therefore if $i> 0$ we have $d(t_i, t_j) \le 2D' +K$.

Now let $x_0, \dots, x_n$ be a $K$-connected path in $T$. We assume $x_0 \in T_a$ and $x_n \in T_b$ with $0<a<b$. The other cases are similar (and simpler). Let $x_{i_a}$ be the first point for which $x_{i_a+1} \not\in  T_a$ and let $x_{i_b}$ be the last point for which $x_{i_b-1} \not\in T_b$.

Form a new sequence $y_{i_a}, \dots, y_{i_b}$ where $y_i = x_i$ if $x_i \in T_0$ and $y_i = t_j$ if $x_i \in T_j$ with $j \neq 0$. It is possible that $x_i$ will be contained in more than one $T_j$. In this case we choose $T_j$ arbitrarily. We claim that $y_0 = t_a, y_1, \dots, y_n = t_b$ is a $(2D'+K)$-connected path. There are four cases:
\begin{itemize}
\item If $x_i, x_{i+1} \in T_j$ with $j \neq 0$ then $d(y_i, y_{i+1}) = d(t_j, t_j) = 0$. 
\item If $x_i, x_{i+1} \in T_0$ then $d(y_i,y_{i+1})= d(x_i, x_{i+1}) \le K$.
\item If $x_i \in T_j$ and $x_{i+1} \in T_k$ with $j,k\neq 0$ and $j\neq k$ then $d(y_i, y_{i+1}) = d(t_j, t_k) \le 2D'+K$.
\item If $x_i \in T_0$ and $x_{i+1} \in T_j$ with $j\neq 0$ then $d(y_i, y_{i+1}) = d(x_i,t_j) \le 2D'$. This similarly holds if $x_i \in T'_j$ and $x_{i+1} \in T_0$.
\end{itemize}

The vertices in the path $\{y_{i_a},\ldots,y_{i_b}\}$ are either in the path $\{x_0,\ldots,x_n\}$ or in $\{t_j\}_j$. If $t_j$ is in $\{y_{i_a},\ldots,y_{i_b}\}$ then there is $x_i \in T_j$ with $x_{i+1}$ or $x_{i-1}$ not in $T_j$. Therefore $d(x_i, t_j) \le 2D'$ and $\{y_{i_a},\ldots,y_{i_b}\}$ is therefore contained in the $2D'$-neighborhood of $\{x_0,\ldots,x_n\}$.

The sets $T_i$ are $Q$-quasi-convex, uniform quality quasi-trees in $X$. Therefore, by Lemma \ref{bottleneck_converse}, there is $\Delta>0$ such that they satisfy the $(\Delta,2D'+K)$-bottleneck condition and
$[x_0, x_{i_a}] \subset B_\Delta(\{x_0, \dots, x_{i_a}\})$,
$ [y_{i_a}, y_{i_b}] \subset  B_\Delta(\{y_{i_a}, \dots, y_{i_b}\})$, and 
$ [x_{i_b}, x_n] \subset  B_\Delta(\{x_{i_b}, \dots, x_n\})$.

 Thus $B_{\Delta +2D'}(\{x_0, \dots, x_n\})$ contains the union of geodesics $[x_0, x_{i_a}] \cup [x_{i_a}, y_{i_a}] \cup [y_{i_a}, y_{i_b}] \cup [y_{i_b}, x_{i_b}] \cup [x_{i_b}, x_n]$. By $\delta$-hyperbolicity, a $4\delta$-neighborhood of this last union contains $[x_0, x_n]$ so
 $$[x_0, x_n] \subset B_{\Delta + 2D' + 4\delta}(\{x_0, \dots, x_n\})$$ and $T$ satisfies the $(\Delta + 2D'+4\delta, K)$-bottleneck condition and is quasi-tree by Corollary \ref{geodesic_bottleneck}.
 
In general, we do not have a uniform bound on the diameter of $I_{Q+2\delta}(T_0, T_i)$. However, we can reduce to this case by applying Lemma \ref{subtrees} to each pair $T_0$ and $T_i$. Each time we do this $T_0$ is enlarged to a subset $T_{0,i}$ (this is $T'_0$ in the lemma) with $T_0 \subset T_{0,i} \subset B_{C'}(T_0)$. Hence if we replace $T_0$ with the union $T'_0 =  \cup T_{0,i}$ we get a $T_0 \subset T'_{0} \subset B_{C'}(T_0)$ so $T'_0$ is a quasi-tree with uniform constants. At each application of the lemma the $T_i$ become quasi-trees $T^j_i$. Then $T'_0$ and the $T^j_i$ have the property that $I_{Q^*+2\delta}(T'_0, T^j_i)$ has uniformly bounded diameter and we have reduced to this case.

Now we complete the induction step. Assume the theorem holds for $N-1$. By assumption, for $i_1, \ldots, i_N$ distinct, the diameter of $I_{Q+2\delta}(T_{i_1}, \dots, T_{i_N})$ is $<D$.
Furthermore, if $I_{Q+2\delta}(T_{i_1}, \dots, T_{i_{N-1}})$ is non-empty then, by Lemma \ref{fellow travel} part \eqref{intersection quasiconvex}, it is a $(Q+4\delta)$-quasi-convex subset and hence a quasi-tree. Denote the set $I_{Q+2\delta}(T_{i_1},\ldots,T_{i_{N-1}})$ by $J_{{i_1}, \dots, {i_{N-1}}}$. Let $J = J_{{i_1}, \dots, {i_{N-1}}}$ and $J'=J_{{i'_1}, \dots, {i'_{N-1}}}$ be two distinct sets of this type. We can assume that $i'_1 \not\in \{i_1, \dots, i_{N-1}\}$. Then $I_{Q+6\delta}(J, J') \subset I_{2Q+ 8\delta}(T_{i_1}, \dots, T_{i_{N-1}}, T_{i'_1})$ and therefore, by Lemma \ref{fellow travel} part \eqref{larger}, the diameter of $I_{Q+2\delta}(J, J')$ is $< D'$ where $D' = D+4Q+20\delta$. By the base case, the union of $T_0$ and the trees $J_{{i_1}, \dots, {i_{N-1}}}$ is a quasi-tree. Label this union $J_0$.

Next apply Lemma \ref{subtrees} to $J_0$ and each $T_i$ to obtain quasi-trees $T'_0$ and $T^j_i$. 
We'll show that $T'_0$ and the sets $T^j_i$ satisfy the theorem conditions for $N-1$. First, we observe that $T'_0$ and the sets $T^j_i$ are $Q^*$-quasi-convex, by Lemma \ref{subtrees}, where $Q^*$ is bounded in terms of our earlier constants.
 
 We claim that the diameter of $I_{Q^*+2\delta}(T^{j_1}_{i_1}, \dots, T^{j_{N-1}}_{i_{N-1}})$ is uniformly bounded. If $i_k = i_\ell$ for some $k \neq \ell$ then $I_{Q^*+2\delta}(T_{i_k}^{j_k}, T_{i_\ell}^{j_\ell})$ is bounded in terms of the constant from Lemma \ref{subtrees}, so that $I_{Q^*+2\delta}(T^{j_1}_{i_1}, \dots, T^{j_{N-1}}_{i_{N-1}})$ is also bounded as a subset. Therefore we can assume that all the $i_k$ are distinct. If the
 diameter of $I_{Q^*+2\delta}(T^{j_1}_{i_1}, \dots, T^{j_{N-1}}_{i_{N-1}})$ is at most $4Q^* + 16\delta$ then there is nothing to show. If not, then the diameter of $I_{Q^*+2\delta}(T_{i_1}, \dots, T_{i_{N-1}})$ is greater than $4Q^* + 16\delta$ and
\begin{eqnarray*}
I_{Q^*+2\delta}\left(T^{j_1}_{i_1}, \dots, T^{j_{N-1}}_{i_{N-1}}\right) &\subset & I_{Q^*+2\delta}(T_{i_1}, \dots, T_{i_{N-1}}) \\
& \subset & B_{Q^* + 4\delta}(J_{i_1, \dots, i_{N-1}}) \\ %I_{Q+2\delta}(T^{j_1}_{i_1}, \dots, T^{j_{N-1}}_{i_{N-1}}) \\
& \subset & B_{Q^*+4\delta}(T'_0)
\end{eqnarray*}
where the middle inclusion follows from Lemma \ref{fellow travel} part \eqref{larger}.
It follows that
\begin{eqnarray*}
I_{Q^*+2\delta}\left(T^{j_1}_{i_1}, \dots, T^{j_{N-1}}_{i_{N-1}}\right)& = &I_{Q^*+2\delta}\left(T^{j_1}_{i_1}, \dots, T^{j_{N-1}}_{i_{N-1}}\right) \cap B_{Q^*+4\delta}(T'_0)\\ & = & I_{Q^*+4\delta}\left(T^{j_1}_{i_1}, \dots, T^{j_{N-1}}_{i_{N-1}}, T'_0\right) \\
&\subset& I_{Q^*+4\delta}\left(T^{j_1}_{i_1}, T'_0\right) \cup \cdots \cup I_{Q^*+4\delta}\left(T^{j_{N-1}}_{i_{N-1}}, T'_0\right).
\end{eqnarray*}
As the final sets on the right are uniformly bounded and $I_{Q^*+2\delta}(T^{j_1}_{i_1}, \dots, T^{j_{N-1}}_{i_{N-1}})$ is coarsely connected (with controlled bounds) this gives a uniform bound on $I_{Q^*+2\delta}(T^{j_1}_{i_1}, \dots, T^{j_{N-1}}_{i_{N-1}})$. This completes the induction step of the proof.
\end{proof}

\end{document}
